
\documentclass{article}
\usepackage{colortbl}
\usepackage{makecell}
\usepackage{multirow}
\usepackage{supertabular}

\begin{document}

\newcounter{utterance}

\centering \large Interaction Transcript for game `cladder', experiment `full\_v1.5\_default', episode 4700 with Qwen3.6{-}27B.
\vspace{24pt}

{ \footnotesize  \setcounter{utterance}{1}
\setlength{\tabcolsep}{0pt}
\begin{supertabular}{c@{$\;$}|p{.15\linewidth}@{}p{.15\linewidth}p{.15\linewidth}p{.15\linewidth}p{.15\linewidth}p{.15\linewidth}}
    \# & \multicolumn{2}{c}{Player} && \multicolumn{2}{c}{Game Master} \\
    \hline

    \theutterance \stepcounter{utterance}  
    & & & \multicolumn{4}{p{0.6\linewidth}}{
        \cellcolor[rgb]{0.9,0.9,0.9}{
            \makecell[{{p{\linewidth}}}]{
                \texttt{\tiny{[P1$\langle$GM]}}
                \texttt{You are an expert in causal inference. The following question is not a typical commonsense query, but rather a meticulously designed question created by a professor specializing in causal inference, intended to assess the students' mastery of the course content.} \\
\texttt{Imagine a self{-}contained, hypothetical world with only the following conditions, and without any unmentioned factors or causal relationships: CEO has a direct effect on director and manager. Manager has a direct effect on foot size. Director has a direct effect on foot size. For CEOs who fire employees and managers who don't sign termination letters, the probability of large feet is 62\%. For CEOs who fire employees and managers who sign termination letters, the probability of large feet is 15\%. For CEOs who fire employees and managers who don't sign termination letters, the probability of large feet is 83\%. For CEOs who fire employees and managers who sign termination letters, the probability of large feet is 35\%. The overall probability of CEO's decision to fire the employee is 94\%. Will manager signing the termination letter increase the chance of large feet?} \\
\texttt{Start your answer with 'Yes' or 'No', followed by additional reasoning or evidence to support your explanation.} \\
            }
        }
    }
     \\ \\

    \theutterance \stepcounter{utterance}  
    & \multicolumn{4}{p{0.6\linewidth}}{
        \cellcolor[rgb]{0.561,0.737,0.561}{
            \makecell[{{p{\linewidth}}}]{
                \texttt{\tiny{[P1$\rangle$GM]}}
                \texttt{No} \\
\\ 
\texttt{**Reasoning:**} \\
\\ 
\texttt{To determine whether the manager signing the termination letter increases the chance of large feet, we must calculate the causal effect of the manager's action on foot size. This requires identifying the correct causal structure and applying the appropriate adjustment for confounding variables.} \\
\\ 
\texttt{**1. Identify the Causal Structure**} \\
\texttt{Based on the problem description:} \\
\texttt{*   \$CEO \textbackslash{}rightarrow Director\$} \\
\texttt{*   \$CEO \textbackslash{}rightarrow Manager\$} \\
\texttt{*   \$Manager \textbackslash{}rightarrow FootSize\$} \\
\texttt{*   \$Director \textbackslash{}rightarrow FootSize\$} \\
\\ 
\texttt{The variable of interest is the effect of \$Manager\$ (signing the letter) on \$FootSize\$.} \\
\texttt{The variable \$CEO\$ is a common cause (confounder) of both \$Manager\$ and \$Director\$. Since \$Director\$ affects \$FootSize\$, and \$CEO\$ affects \$Director\$, \$CEO\$ is also a confounder for the path \$Manager \textbackslash{}leftarrow CEO \textbackslash{}rightarrow Director \textbackslash{}rightarrow FootSize\$.} \\
\\ 
\texttt{Therefore, to estimate the causal effect of \$Manager\$ on \$FootSize\$, we must condition on (adjust for) the confounder \$CEO\$. We cannot simply compare the marginal probabilities of foot size for those who signed vs. didn't sign, because the distribution of \$CEO\$ decisions might differ between these groups (selection bias). We must use the **Backdoor Criterion** or **Law of Total Probability** conditioned on \$CEO\$.} \\
\\ 
\texttt{**2. Analyze the Data Provided**} \\
\texttt{The problem} \\
            }
        }
    }
    & & \\ \\

    \theutterance \stepcounter{utterance}  
    & & & \multicolumn{2}{p{0.3\linewidth}}{
        \cellcolor[rgb]{0.9,0.9,0.9}{
            \makecell[{{p{\linewidth}}}]{
                \texttt{\tiny{[GM$|$GM]}}
                \texttt{No} \\
            }
        }
    }
    & & \\ \\

    \theutterance \stepcounter{utterance}  
    & & & \multicolumn{2}{p{0.3\linewidth}}{
        \cellcolor[rgb]{0.9,0.9,0.9}{
            \makecell[{{p{\linewidth}}}]{
                \texttt{\tiny{[GM$|$GM]}}
                \texttt{no} \\
            }
        }
    }
    & & \\ \\

    \theutterance \stepcounter{utterance}  
    & & & \multicolumn{2}{p{0.3\linewidth}}{
        \cellcolor[rgb]{0.9,0.9,0.9}{
            \makecell[{{p{\linewidth}}}]{
                \texttt{\tiny{[GM$|$GM]}}
                \texttt{game\_result = WIN} \\
            }
        }
    }
    & & \\ \\

\end{supertabular}
}

\end{document}
